\section{机器人数据瓶颈：真实世界太复杂，数据太贵}

前面说模型还没有完全独立，本章进入机器人最硬的问题：数据。谭杰认为机器人最大问题是数据。机器人处在复杂的 unstructured environment 中，任何事情都可能发生；动作轨迹昂贵，长尾失败难覆盖，硬件本体不同还会导致数据难以共享。

\begin{figure}[H]
\centering
\includegraphics[width=0.96\textwidth]{robotics-data-bottleneck.png}
\caption{机器人数据瓶颈：非结构环境、长尾失败和动作反馈让数据特别贵。自制概念图，依据 00:23:44--00:27:52 对谈内容整理。}
\end{figure}

\begin{knowledgebox}{读图：机器人数据不是普通视频数据}
真实环境不可控，动作轨迹采集昂贵，失败样本难覆盖，跨本体有硬件差异，精细操作成功率低。高质量数据因此成为机器人从 demo 到泛化的核心瓶颈。
\end{knowledgebox}

\subsection{为什么几百小时不够验证 scaling}

本节解释为什么机器人数据问题比普通视频或文本数据更难。训练机器人不是只收集“看起来相关”的视频，而是要收集带动作、结果、失败和恢复路径的数据。

谭杰提到，在机器人上如果只有几百小时数据，可能无法验证 scaling；你可能需要几万小时数据，才能看到训练 recipe 是否真的 work。这个要求对创业团队和投资人都很残酷，因为初期小钱小数据很难证明长期路线，而机器人又特别需要长期信仰和大投入。

\begin{warningbox}{小数据 demo 的误导}
几十小时或几百小时数据可能做出漂亮 demo，但不足以证明泛化。机器人真正需要的是跨场景、跨物体、跨本体、跨失败模式的高质量数据。
\end{warningbox}

\subsection{特斯拉路线为什么难复制到通用机器人}

特斯拉汽车有用，所以每天产生真实驾驶数据，数据飞轮能转起来。许多机器人硬件虽然能采数据，但还没有达到“本身有用”的阈值，因此数据飞轮转不起来。这是为什么不能简单说“中国有硬件所以能复制特斯拉路线”：硬件必须先进入真实使用，数据才有持续回流。

\subsection{机器人像几岁小孩：能力不均衡}

本节补充一个很容易被忽视的判断：机器人能力不是均匀长大的。谭杰用小孩做类比，认为机器人的 locomotion 已经很强，甚至某些运动能力超过成年人；但 manipulation，尤其是灵巧手操作，可能还像两三岁小孩，能大概理解指令、尝试几次，但抓取不稳，精细动作更难。

\begin{knowledgebox}{能力不均衡：Locomotion vs Manipulation}
\begin{tabularx}{\textwidth}{p{0.24\textwidth}p{0.34\textwidth}X}
\toprule
能力 & 当前状态 & 为什么不同步 \\
\midrule
Locomotion & 跑、跳、平衡和步态已经进展很快 & 强化学习和 Sim2Real 在过去五年基本解决了许多运动控制问题。 \\
普通抓取 & 能完成简单抓放和跟随指令 & 需要视觉、目标理解和粗控制，但容错相对高。 \\
灵巧手操作 & 仍然很难，很多任务不稳定 & 多指协同、触觉、接触力和材质反馈都很复杂。 \\
认知规划 & 大模型带来常识和任务拆解 & 但计划仍要和真实动作成功率耦合。 \\
\bottomrule
\end{tabularx}
\end{knowledgebox}

\subsection{本章小结}

机器人领域的第一性问题仍是高质量数据。没有足够真实、足够多样、足够可验证的数据，就无法判断模型是否能 scale，也无法从 demo 走向可用产品。

